%! Random Sampling Methods — Unit 3, Lesson 3.2 — Warm-Up
% ONE PAGE, blank and key. Four short items; do not add a fifth.
% GRADUAL RELEASE: the warm-up activates prior knowledge before the guided notes.
%   Item 1 — spiral to Lesson 3.1: the population and the sample in the Lakeview survey.
%   Item 2 — names in a hat: does everyone have the same chance? Seeds the SRS row (row 1).
%   Item 3 — the hat could, by chance, draw few seniors; seniors work more. Seeds STRATA (row 2/3).
%   Item 4 — surveying two whole advisory groups is easier. Seeds CLUSTERS and the "cheaper"
%            advantage (row 3).
% Every display is given; nothing is sketched.
\documentclass[12pt]{article}
\usepackage{apstats-article}
\usepackage{apstats-boxes}

\begin{document}

\pageheader{Unit 3, Lesson 3.2}{Warm-Up}
% NAMESTRIP: no \namedateperiod here — the name row lives on the cover only.

\vspace{0.22in}

The principal of Lakeview High (1{,}200 students, 300 in each of grades 9--12) wants to know how
many hours per week her students work at a paid job. She plans to survey 80 students.

\vspace{0.16in}

\begin{enumerate}[leftmargin=1.9em, itemsep=0.34in]

  \item \textbf{From Lesson 3.1.} Name the population and the sample.

        \vspace{12pt}
        Population: \blank{4.6cm} \qquad Sample: \blank{4.6cm}

  \item She puts all 1{,}200 names in a hat and draws 80. Does every student have the same
        chance of being picked? Explain.

        \vspace{12pt}
        \writeline\\[18pt]\writeline

  \item Seniors usually work many more hours than freshmen. Just by chance, the hat could give
        her very few seniors. What would that do to her estimate? Suggest a way to choose the 80
        that makes sure every grade is in the sample.

        \vspace{12pt}
        \writeline\\[18pt]\writeline\\[18pt]\writeline

  \item Lakeview's 48 advisory groups each have students from all four grades. Why might it be
        \emph{easier} to survey every student in 2 advisory groups than to find 80 students
        scattered across the building?

        \vspace{12pt}
        \writeline\\[18pt]\writeline

\end{enumerate}

\end{document}
